% Authors: Teddy Wachtler
%%%%%%%%%%%%%%%%%%%%%%%%%%%%%%%%%%%%%%%%%%%%%%%

\documentclass{article}
\usepackage[dvips]{graphicx}
\usepackage{a4wide}
\usepackage{amsmath}
\usepackage{euscript}
\usepackage{amssymb}
\usepackage{amsthm}
\usepackage{amsopn}

\theoremstyle{definition}
\newtheorem*{definition}{Definition}
\newtheorem{theorem}{Theorem}

\renewcommand{\qed}{\blacksquare}
\newcommand{\powset}{{\cal P}}
\newcommand{\img}{_{*}}
\newcommand{\pimg}{^{*}}

\newcommand{\vv}{\ensuremath{\vec{v}}}
\newcommand{\vu}{\ensuremath{\vec{u}}}
\newcommand{\vw}{\ensuremath{\vec{w}}}
\newcommand{\vx}{\ensuremath{\vec{x}}}
\newcommand{\vy}{\ensuremath{\vec{y}}}
\newcommand{\vb}{\ensuremath{\vec{b}}}
\newcommand{\vo}{\ensuremath{\vec{0}}}
\newcommand{\va}{\ensuremath{\vec{a}}}
\newcommand{\ve}{\ensuremath{\vec{e}}}

\newcommand{\R}{\mathbb{R}}
\newcommand{\Z}{\mathbb{Z}}
\newcommand{\C}{\mathbb{C}}
\newcommand{\N}{\mathbb{N}}
\newcommand{\Q}{\mathbb{Q}}

%%%%%%%%%%%%%%%%%%%%%%%%%%%%%%%%%%%%%%%%%%%%%%%

\begin{document}
	\begin{center}
		\Large{Math 251W: Foundations of Advanced Mathematics}
		\vspace{0.1cm}

		\normalsize Portfolio Assignment 4: \S 4.3 \& 4.4
		\vspace{0.1cm}

		\hfill Name: Teddy Wachtler
		\vspace{0.1cm}

		\hrule
		\vspace{0.5cm}
	\end{center}

%%%%%%%%%%%%%%%%%%%%%%%%%%%%%%%%%%%%%%%%%%%%%%%

	\noindent Problem 4.4.14
	\vspace{0.2cm}

	\begin{enumerate}
		\item[3]
		\underline{Proposition:}
		Given functions $f : A \rightarrow B$, and $g : B \rightarrow C$, if $g \circ f$ is bijective, then $f$ must be injective and $g$ must be surjective.

	\end{enumerate}

	\vspace{0.2cm}
	\fbox{Proof} (Direct)
	\vspace{0.2cm}

	Let $A, B$, and $C$ be arbitrary sets.
	Let $f$ and $g$ be arbitrary functions defined by $f : A \rightarrow B$ and $g : B \rightarrow C$.
	Suppose $g \circ f$ is bijective.
	By the definition of bijective functions, we see that $g \circ f$ is injective and surjective.

	Let $a_1$ and $a_2$ be arbitrary elements in $A$ so that $f(a_1) = f(a_2)$
	By substitution, we see that $g(f(a_1)) = g(f(a_2))$.
	By the definition of function composition, we see that $(g \circ f)(a_1) = (g \circ f)(a_2)$.
	By the definition of injective functions, we see that $a_1 = a_2$.
	By the definition of injective functions, since $a_1$ and $a_2$ were arbitrary elements of $A$, and $f(a_1) = f(a_2)$ implies $a_1 = a_2$, we see that $f$ is injective.

	Let $c$ be an arbitrary element in $C$.
	By the definition of surjective functions, since $g \circ f$ is surjective, we see that there exists $a \in A$ so that $(g \circ f)(a) = c$.
	By the definition of function composition, we see that $g(f(a)) = c$.
	Since $f : A \rightarrow B$, we see that there exists $b \in B$ so that $f(a) = b$.
	By substitution, we see that $g(b) = c$.
	Therefore, by the definition of surjective functions, since $c$ was an arbitrary element in $C$, and there exists an element $b \in B$ so that $g(b) = c$, we see that $g$ is surjective.

	Therefore, we see that for all sets $A, B$, and $C$, and for all functions $f : A \rightarrow B$, and $g : B \rightarrow C$, if $g \circ f$ is bijective, then $f$ must be injective and $g$ must be surjective.
	$\qed$.

	\vspace{0.5cm}

%%%%%%%%%%%%%%%%%%%%%%%%%%%%%%%%%%%%%%%%%%%%%%%

	\newpage

	\noindent Challenge Problem
	\vspace{0.2cm}

	\underline{Proposition:}
	Let $f: A \rightarrow B$ be a function.
	We can define a new function $F: \mathcal{P}(A) \rightarrow \mathcal{P}(B)$ by the image map:\[ F(S) = f_{*}(S) = \{ f(x)  \mid x \in S \} \]

	\begin{enumerate}
		\item
		Prove that if $f$ is surjective, then $F$ is surjective.

		\vspace{0.2cm}
		\fbox{Proof} (Direct)
		\vspace{0.2cm}

		Let $A$ and $B$ be arbitrary sets.
		Let $f$ be an arbitrary function defined by $f : A \rightarrow B$.
		Let $F$ be an arbitrary function defined by $F : \powset (A) \rightarrow \powset (B)$ and $F(S) = f \img (S) = \{ f(x) : x \in S \}$.
		Suppose $f$ is surjective.
		Let $Y$ be an arbitrary element of $\powset(B)$.
		Let $X$ be the set defined by $X = \{x \in A : f(x) \in Y\}$.

		Let $y_1$ be an arbitrary element of $f \img (X)$.
		By the definition of the image, we see that there exists a $x_1 \in X$ so that $f(x_1) = y_1$.
		By the definition of $X$, since $x_1 \in X$, we see that $f(x_1) \in Y$.
		By substitution, we see that $y_1 \in Y$.
		Therefore, by the definition of subsets, since $y_1 \in f \img (X)$ implies $y_1 \in Y$, we see that $f \img (X) \subseteq Y$.

		Let $y_2$ be an arbitrary element of $Y$.
		By the definition of the power set, we see that $Y \subseteq B$.
		By the definition of subsets, since $y_2 \in Y$, we see that $y_2 \in B$.
		By the definition of surjective functions, since $f$ is surjective, we see that there exists $x_2 \in A$ so that $f(x_2) = y_2$.
		By the definition of $X$, since $f(x_2) \in Y$, we see that $x_2 \in X$.
		By the definition of the image, we see that $y_2 \in f \img (X)$
		Therefore, by the definition of subsets, since $y_2 \in Y$ implies $y_2 \in f \img (X)$, we see that $Y \subseteq f \img (X)$.

		Therefore, by the definition of set equality, since $Y \subseteq f \img (X)$ and $f \img (X) \subseteq Y$, we see that $Y = f \img (X)$.
		By the definition of $X$, we see that $X \subseteq A$.
		By the definition of the power set, we see that $X \in \powset (A)$.
		By the definition of $F$, we see that $F(X) = Y$.
		By the definition of surjective functions, since $Y$ was an arbitrary element of $\powset (B)$ and there exists an element $X$ so that $F(X) = Y$, we see that $F$ is surjective.
		Therefore, for all sets $A$ and $B$, and for all functions $f$ defined by $f : A \rightarrow B$ and functions $F$ defined by $F : \powset (A) \rightarrow \powset (B)$ and $F(S) = f \img (S) = \{ f(x) : x \in S \}$, we see that if $f$ is surjective, then $F$ is surjective.
		$\qed$.

		\vspace{0.2cm}
		\item
		Is it true that if $f$ is injective, then $F$ is injective?
		Provide a Proof or counterexample.

		\vspace{0.2cm}
		\fbox{Proof} (Direct)
		\vspace{0.2cm}

		Let $A$ and $B$ be arbitrary sets.
		Let $f$ be an arbitrary function defined by $f : A \rightarrow B$.
		Let $F$ be an arbitrary function defined by $F : \powset (A) \rightarrow \powset (B)$ and $F(S) = f \img (S) = \{ f(x) : x \in S \}$.
		Suppose $f$ is injective.
		Let $A_1$ and $A_2$ be arbitrary elements of $\powset (A)$ so that $F(A_1) = F(A_2)$.
		By the definition of $F$, we see that $f \img (A_1) = f \img (A_2)$.
		Let $a_1$ be an arbitrary element in $A_1$.
		By the definition of $f$, there exists a particular element $b \in B$ so that $f(a_1) = b$.
		By the definition of images, we see that $b \in f \img (A_1)$.
		By the definition of set equality, since $f \img (A_1) = f \img (A_2)$, we see that $b \in f \img (A_2)$.
		By the definition of the image, we see that there exists some element $a_2 \in A_2$ so that $f(a_2) = b$.
		By the definition of injective functions, since $f$ is injective, we see that $a_1 = a_2$.
		By substitution, we see that $a_1 \in A_2$.
		By the definition of subsets, since $a_1 \in A_1$ implies $a_1 \in A_2$, we see that $A_1 \subseteq A_2$.
		By a similar argument, we see that $A_2 \subseteq A_1$.
		By the definition of set equality, we see that $A_1 = A_2$.
		By the definition of injective functions, since $A_1$ and $A_2$ were arbitrary elements of $\powset (A)$ and $F(A_1) = F(A_2)$ implies $A_1 = A_2$, we see that $F$ is injective.
		Therefore, for all sets $A$ and $B$, and for all functions $f$ defined by $f : A \rightarrow B$ and functions $F$ defined by $F : \powset (A) \rightarrow \powset (B)$ and $F(S) = f \img (S) = \{ f(x) : x \in S \}$, we see that if $f$ is injective, then $F$ is injective.
		$\qed$.

		\vspace{0.2cm}
	\end{enumerate}

%%%%%%%%%%%%%%%%%%%%%%%%%%%%%%%%%%%%%%%%%%%%%%%

\end{document}